% Generated draft: requires mathematical and source review.
\documentclass[11pt]{article}
\usepackage[margin=25mm]{geometry}
\usepackage{fontspec}
\setmainfont{Latin Modern Roman}
\usepackage{amsmath,amssymb,mathtools,unicode-math}
\setmathfont{Latin Modern Math}
\usepackage{enumitem,needspace,xurl,hyperref}
\setlength{\parindent}{0pt}
\setlength{\parskip}{5pt}
\setlength{\emergencystretch}{3em}
\newcommand{\N}{\mathbb{N}}
\newcommand{\Z}{\mathbb{Z}}
\newcommand{\Q}{\mathbb{Q}}
\newcommand{\R}{\mathbb{R}}
\newcommand{\C}{\mathbb{C}}
\newcommand{\F}{\mathbb{F}}
\newcommand{\A}{\mathbb{A}}
\newcommand{\PP}{\mathbb{P}}
\newcommand{\E}{\mathbb{E}}
\newcommand{\eps}{\varepsilon}
\newcommand{\dd}{\,\mathrm{d}}
\begin{document}
\section*{UnsolvedMath ID: 30003873}
\textbf{Source problem code:} catalogue entry
\subsection{Definitions and mathematical statement}
The base field and the precise model of the universal polynomial cover are not specified in the supplied record. The definitions below give a conventional algebraic interpretation; identifying it with the intended formulation requires source review. In particular, the statement below does not assert a conjecture uniformly over all fields.

Let \(k\) denote the intended, presently unspecified base field. A conventional parameter space for monic degree-\(n\) polynomials is
\[
P_n=\mathbb A_k^n,\qquad f_{\mathbf a}(T)=T^n+a_1T^{n-1}+\cdots+a_n.
\]
The corresponding single-root cover is
\[
\widetilde P_n=\{(\mathbf a,t)\in P_n\times\mathbb A_k^1:f_{\mathbf a}(t)=0\},\qquad
\pi_n(\mathbf a,t)=\mathbf a.
\]
This model includes polynomials with repeated roots; its resolvent degree is defined at the generic point. Write
\[
K=k(a_1,\ldots,a_n),\qquad L=K[T]/(T^n+a_1T^{n-1}+\cdots+a_n),
\]
where the coefficients are algebraically independent over \(k\). The generic polynomial is irreducible, so \(L\) is a field.

For a finite extension \(L/K\) of finitely generated fields over \(k\), define its algebraic resolvent degree \(\operatorname{RD}_k(L/K)\) to be the least integer \(d\geq0\) for which, inside a common algebraic closure of \(K\), there is a finite tower
\[
K=K_0\subseteq K_1\subseteq\cdots\subseteq K_r,\qquad L\subseteq K_r,
\]
such that for every \(i\) there are finitely generated fields \(F_i,E_i\) over \(k\) satisfying
\[
F_i\subseteq K_{i-1},\qquad E_i\subseteq K_i,\qquad
[E_i:F_i]<\infty,\qquad K_i=K_{i-1}E_i,\qquad
\operatorname{trdeg}_k(F_i)\leq d.
\]
For the single-root model, put \(\operatorname{RD}_k(\pi_n)=\operatorname{RD}_k(L/K)\). This specifies the convention used here; the supplied record does not itself define resolvent degree or identify the base field.

With those qualifications, the conjectural target attributed to Hilbert in the supplied record is
\[
\text{for every integer }n\geq6,\qquad
\operatorname{RD}(\widetilde P_n\to P_n)\geq2,
\]
together with the more specific conjectural equalities
\[
\begin{aligned}
\operatorname{RD}(\widetilde P_6\to P_6)&=2,&
\operatorname{RD}(\widetilde P_7\to P_7)&=3,\\
\operatorname{RD}(\widetilde P_8\to P_8)&=4,&
\operatorname{RD}(\widetilde P_9\to P_9)&=4.
\end{aligned}
\]
Here the unsubscripted notation preserves the source statement; under the interpretation above it means \(\operatorname{RD}_k\) for the intended fixed field. The source does not specify whether \(\widetilde P_n\) selects one root or parametrizes an ordered collection of all roots. That identification remains unresolved. These assertions are supplied as conjectures, not as established mathematical facts.
\subsection{Short English statement}
The supplied conjecture says that solving the generic degree-\(n\) polynomial requires algebraic functions of at least two parameters for every integer \(n\geq6\). It further predicts resolvent degrees \(2,3,4,4\) in degrees \(6,7,8,9\), respectively. The intended base field and universal-cover convention require confirmation. The supplied open-status assessment is dated curation and has not been independently verified.
\subsection{Sources}
\begin{itemize}
\item Local UnsolvedMath registry, supplied complete record for permanent ID \(30003873\), title ``Resolvent Degrees of Universal Polynomials,'' exact source problem code \texttt{catalogue entry}; collection: Oberwolfach Reports Open Problems. The record attributes the conjecture to Hilbert but supplies no separately named proposer.
\item Supplied report citation: \emph{Algebraische Zahlentheorie (2019)}, Oberwolfach Reports, DOI \texttt{10.4171/owr/2018/30}, pp.~1785--1844. \url{https://doi.org/10.4171/owr/2018/30}. The background describes Oberwolfach Reports as documenting mathematical workshops and their open-problem discussions. The record separately gives proposed year \(2018\) and the tag \texttt{owr-year:2018}; the relationship to the citation's \(2019\) date is not explained.
\item Edens--Reichstein, ``Hilbert's 13th problem,'' as supplied in the literature sources. \url{https://arxiv.org/abs/2406.15954}.
\item Farb--Wolfson, ``resolvent degree,'' as supplied in the literature sources. \url{https://arxiv.org/abs/1803.04063}.
\item Supplied literature assessment, checked on \texttt{2026-08-21}: status ``open.'' The assessment states that the lower bounds remain conjectural and that the literature supplies upper bounds and reformulations but not the claimed general lower bound. This is supplied dated curation derived from the OWR record, not independently verified source consultation or a new literature assessment; time-sensitive status and cited claims require independent verification.
\end{itemize}
\end{document}
